\documentclass[11pt, a4paper]{article}

% bfw-style.tex — gemeinsame Style-Definitionen für alle Handouts/Aufgabenhefte
% Voraussetzung: Kompilation mit xelatex (wegen fontspec)

\usepackage[ngerman]{babel}
\usepackage{geometry}
\geometry{a4paper, top=2.2cm, bottom=2.2cm, left=2.3cm, right=2.3cm}

\usepackage{fontspec}
% Fallback-Kette: Source Sans Pro -> Calibri -> Segoe UI -> Latin Modern Sans
% (Latin Modern ist immer mit MiKTeX vorhanden)
\IfFontExistsTF{Source Sans Pro}{%
  \setmainfont{Source Sans Pro}[Ligatures=TeX]%
}{\IfFontExistsTF{Calibri}{%
  \setmainfont{Calibri}[Ligatures=TeX]%
}{\IfFontExistsTF{Segoe UI}{%
  \setmainfont{Segoe UI}[Ligatures=TeX]%
}{%
  \setmainfont{Latin Modern Sans}[Ligatures=TeX]%
}}}
\IfFontExistsTF{Source Code Pro}{%
  \setmonofont{Source Code Pro}[Scale=0.92]%
}{\IfFontExistsTF{Consolas}{%
  \setmonofont{Consolas}[Scale=0.92]%
}{%
  \setmonofont{Latin Modern Mono}[Scale=0.92]%
}}

\usepackage{xcolor}
% Farbpalette: hell, freundlich, modern
\definecolor{bfwPrimary}{HTML}{0EA5A4}    % Petrol/Türkis - Primär
\definecolor{bfwPrimaryDark}{HTML}{0F766E} % dunkler Petrol für Headings
\definecolor{bfwAccent}{HTML}{F59E0B}     % Amber - Akzent
\definecolor{bfwSuccess}{HTML}{10B981}    % Grün - Erfolg / Lösung
\definecolor{bfwWarn}{HTML}{F97316}       % Orange - Achtung
\definecolor{bfwInk}{HTML}{0F172A}        % fast Schwarz
\definecolor{bfwInkSoft}{HTML}{475569}    % gedämpftes Grau
\definecolor{bfwBgSoft}{HTML}{F8FAFC}     % off-white
\definecolor{bfwBgTeal}{HTML}{ECFEFF}     % helles Türkis
\definecolor{bfwBgAmber}{HTML}{FEF3C7}    % helles Gelb
\definecolor{bfwBgRose}{HTML}{FFE4E6}     % helles Rosé für Eselsbrücken

\usepackage[colorlinks=true, linkcolor=bfwPrimaryDark, urlcolor=bfwPrimaryDark]{hyperref}
\usepackage{enumitem}
\setlist[itemize]{leftmargin=*, itemsep=2pt, topsep=2pt}
\setlist[enumerate]{leftmargin=*, itemsep=2pt, topsep=2pt}

\usepackage[most]{tcolorbox}
\usepackage{booktabs}
\usepackage{tikz}
\usetikzlibrary{decorations.pathmorphing, positioning, shapes.geometric, arrows.meta}

% Merkkästchen — wichtige Definition / Kernpunkt
\newtcolorbox{merksatz}[1][]{%
  enhanced, breakable,
  colback=bfwBgTeal,
  colframe=bfwPrimary,
  boxrule=0.6pt,
  arc=4pt,
  left=10pt, right=10pt, top=8pt, bottom=8pt,
  fonttitle=\bfseries\color{bfwPrimaryDark},
  title={\faLightbulb\ Merksatz},
  #1
}

% Eselsbrücke — Mnemoniks
\newtcolorbox{eselsbruecke}[1][]{%
  enhanced, breakable,
  colback=bfwBgRose,
  colframe=bfwAccent,
  boxrule=0.6pt,
  arc=4pt,
  left=10pt, right=10pt, top=8pt, bottom=8pt,
  fonttitle=\bfseries\color{bfwWarn},
  title={Eselsbrücke},
  #1
}

% Beispiel-Box
\newtcolorbox{beispiel}[1][]{%
  enhanced, breakable,
  colback=bfwBgSoft,
  colframe=bfwInkSoft,
  boxrule=0.4pt,
  arc=3pt,
  left=10pt, right=10pt, top=8pt, bottom=8pt,
  fonttitle=\bfseries\color{bfwInk},
  title={Beispiel},
  #1
}

% Lösungs-Box (für Aufgabenhefte)
\newtcolorbox{loesung}[1][]{%
  enhanced, breakable,
  colback=bfwBgAmber!40,
  colframe=bfwSuccess,
  boxrule=0.6pt,
  arc=3pt,
  left=10pt, right=10pt, top=8pt, bottom=8pt,
  fonttitle=\bfseries\color{bfwSuccess!70!black},
  title={Lösung},
  #1
}

% Glossar-Eintrag
\newcommand{\glossareintrag}[2]{%
  \par\noindent\textbf{\color{bfwPrimaryDark}#1} \enspace #2\par\smallskip
}

% Section-Stil — etwas farbiger als Standard
\usepackage{titlesec}
\titleformat{\section}
  {\Large\bfseries\color{bfwPrimaryDark}}
  {\thesection}{0.6em}{}
\titleformat{\subsection}
  {\large\bfseries\color{bfwPrimary}}
  {\thesubsection}{0.5em}{}
\titleformat{\subsubsection}
  {\normalsize\bfseries\color{bfwInk}}
  {\thesubsubsection}{0.4em}{}

% TikZ-Stil "skizzenhaft" — etwas weicher als Rough.js, aber stabil
\tikzset{
  roughbox/.style = {
    draw=bfwPrimaryDark, thick, fill=bfwBgTeal,
    rounded corners=4pt, inner sep=6pt
  },
  roughaccent/.style = {
    draw=bfwAccent, thick, fill=bfwBgAmber,
    rounded corners=4pt, inner sep=6pt
  },
  roughline/.style = {
    draw=bfwInkSoft, thick, -{Stealth[length=2.5mm]}
  }
}

% Bullet-Symbole (bewusst kein FontAwesome — vermeidet Font-Installations-Probleme)
\newcommand{\faLightbulb}{$\bullet$}
\newcommand{\faBookmark}{$\bullet$}

% Header/Footer
\usepackage{fancyhdr}
\pagestyle{fancy}
\fancyhf{}
\renewcommand{\headrulewidth}{0pt}
\fancyfoot[L]{\small\color{bfwInkSoft}\thepage}
\fancyfoot[R]{\small\color{bfwInkSoft} BFW M-064 Beschaffung im Büromanagement}


\title{\vspace*{-1.5cm}\Huge\bfseries\color{bfwPrimaryDark}Aufgabenheft Tag~13\\[0.4em]
  \large\color{bfwInkSoft}\textnormal{Ablage \& Ablagesysteme --- Übungen im IHK-Klausurformat}}
\author{\textsf{Büroprozesse gestalten und Arbeitsvorgänge organisieren \textperiodcentered{} M-067}\\
\textsf{\small\copyright\ Can Siebert\,\textperiodcentered\,2026}}
\date{Selbstlernmaterial \textperiodcentered{} 15.07.2026}

\begin{document}
\maketitle
\thispagestyle{empty}

\begin{merksatz}[title={\bfseries So nutzen Sie dieses Aufgabenheft}]
Bearbeiten Sie alle Aufgaben zunächst \emph{ohne} in die Lösungen zu schauen. Schreiben
Sie Ihre Antworten in vollständigen Sätzen --- die IHK-Prüfung verlangt formulierte
Begründungen, nicht bloße Stichworte. Beachten Sie die \emph{Operatoren} (nennen,
erläutern, beurteilen, begründen) --- sie geben den geforderten Antwort-Tiefegrad vor.
Erst danach öffnen Sie den Lösungsteil ab Seite~\pageref{loesungen} und vergleichen.
\end{merksatz}

\tableofcontents
\vspace{1em}
\hrule

\section{Aufgaben}

\subsection*{Aufgabe~1 --- Ablageorte (10 Punkte)}

\noindent\textbf{\color{bfwAccent}YouTube-Vorbereitung:}\enspace \href{https://www.youtube.com/watch?v=t-yOsZ8tDf4}{Ablage \& Registratur --- Grundlagen} (\url{https://www.youtube.com/watch?v=t-yOsZ8tDf4})

\medskip
In der \emph{Duisdorfer BüroKonzept KG} soll die Ablage neu organisiert werden. Die
Geschäftsführung fragt Sie, wo welche Unterlagen aufbewahrt werden sollten.

\begin{enumerate}[label=(\alph*)]
  \item \textbf{Erläutern} Sie den Grundsatz, nach dem über den Ablageort entschieden wird. \hfill (3~P.)
  \item \textbf{Nennen} Sie die drei Ablageorte und \textbf{ordnen} Sie jedem ein Beispiel zu. \hfill (4~P.)
  \item \textbf{Beschreiben} Sie, welche Aufgabe die Zentralregistratur über die reine Aufbewahrung hinaus erfüllt. \hfill (3~P.)
\end{enumerate}

\subsection*{Aufgabe~2 --- Ablagetechnik und Aktenarten (12 Punkte)}

\noindent\textbf{\color{bfwAccent}YouTube-Vorbereitung:}\enspace \href{https://www.youtube.com/watch?v=aSakXzRYYC4}{Loseblatt- vs.\ geheftete Ablage} (\url{https://www.youtube.com/watch?v=aSakXzRYYC4})

\medskip
\begin{enumerate}[label=(\alph*)]
  \item \textbf{Erläutern} Sie je einen Vor- und einen Nachteil der Loseblattablage. \hfill (3~P.)
  \item \textbf{Erläutern} Sie je einen Vor- und einen Nachteil der gehefteten Ablage. \hfill (3~P.)
  \item \textbf{Unterscheiden} Sie Einzelakte und Sammelakte mit je einem Beispiel. \hfill (3~P.)
  \item \textbf{Beurteilen} Sie, für welche Schriftstücke die Loseblattablage besser geeignet ist als die geheftete Ablage. \hfill (3~P.)
\end{enumerate}

\subsection*{Aufgabe~3 --- Arten der Registratur (15 Punkte)}

\noindent\textbf{\color{bfwAccent}YouTube-Vorbereitung:}\enspace \href{https://www.youtube.com/watch?v=7uoooKb-RMU}{Liegende, stehende, hängende Registratur} (\url{https://www.youtube.com/watch?v=7uoooKb-RMU})

\medskip
\begin{enumerate}[label=(\alph*)]
  \item \textbf{Nennen} Sie die drei Aufbewahrungsformen und je eine zugehörige Registraturform. \hfill (3~P.)
  \item \textbf{Ordnen} Sie jeder Registraturform einen passenden Schriftgutbehälter \textbf{zu}. \hfill (4~P.)
  \item \textbf{Erläutern} Sie je einen Vor- und einen Nachteil der Ordnerregistratur. \hfill (4~P.)
  \item \textbf{Vergleichen} Sie Hänge- und Pendelregistratur hinsichtlich Übersicht und Zugriffszeit. \hfill (4~P.)
\end{enumerate}

\subsection*{Aufgabe~4 --- Alphabetische Ordnung nach DIN 5007 (12 Punkte)}

\noindent\textbf{\color{bfwAccent}YouTube-Vorbereitung:}\enspace \href{https://www.youtube.com/watch?v=jZz5wppOvcc}{Alphabetisch sortieren nach DIN 5007} (\url{https://www.youtube.com/watch?v=jZz5wppOvcc})

\medskip
\begin{enumerate}[label=(\alph*)]
  \item \textbf{Sortieren} Sie alphabetisch nach DIN 5007: \emph{Kuhn, Kaiser, Ludwig, Krähmer, Larsson, Kurz}. \hfill (4~P.)
  \item \textbf{Sortieren} Sie alphabetisch nach DIN 5007: \emph{Schmidt, Schmid, Schmitt, Schäfer, Schade, Schulze}. \hfill (4~P.)
  \item \textbf{Erläutern} Sie, wonach sich die Reihenfolge richtet und wie Umlaute sowie „ß” behandelt werden. \hfill (4~P.)
\end{enumerate}

\subsection*{Aufgabe~5 --- Ordnungssystem wählen (12 Punkte)}

\noindent\textbf{\color{bfwAccent}YouTube-Vorbereitung:}\enspace \href{https://www.youtube.com/watch?v=mAffIV-GKt0}{Ordnungssysteme im Büro} (\url{https://www.youtube.com/watch?v=mAffIV-GKt0})

\medskip
\textbf{Ordnen} Sie für jeden Fall ein passendes Ordnungssystem zu \textbf{und begründen} Sie kurz:

\begin{enumerate}[label=\arabic*., leftmargin=*]
  \item Eine Kundenkartei mit 4\,000 Kundennamen.
  \item Eingangs- und Ausgangsrechnungen, abgelegt mit Kürzeln (ER/AR) und Rechnungsnummer.
  \item Bilanzen und Steuerbescheide eines Unternehmens.
  \item Akten eines Außendienstes, geordnet nach Vertriebsgebieten.
\end{enumerate}

\subsection*{Aufgabe~6 --- Ordnungshilfsmittel und Aktenplan (15 Punkte)}

\noindent\textbf{\color{bfwAccent}YouTube-Vorbereitung:}\enspace \href{https://www.youtube.com/watch?v=-O2YBJke99Q}{Aktenplan \& Wiedervorlage} (\url{https://www.youtube.com/watch?v=-O2YBJke99Q})

\medskip
\begin{enumerate}[label=(\alph*)]
  \item \textbf{Erläutern} Sie Zweck und Aufbau der Wiedervorlagemappe. \hfill (4~P.)
  \item \textbf{Nennen} Sie drei weitere Ordnungshilfsmittel mit ihrer Funktion. \hfill (3~P.)
  \item \textbf{Erläutern} Sie, was ein Aktenplan ist und wie er bei einer Dezimalklassifikation aufgebaut ist. \hfill (4~P.)
  \item \textbf{Beurteilen} Sie: Warum ist die günstigste Registraturform nicht automatisch die wirtschaftlichste? \hfill (4~P.)
\end{enumerate}

\vspace{1em}
\noindent\textbf{Punkte gesamt: 76}

\newpage

\section{Lösungen}\label{loesungen}

\subsection*{Lösung Aufgabe~1}
\begin{loesung}
\textbf{(a)} Grundsatz: Je häufiger ein Mitarbeiter auf die Schriftstücke zurückgreifen
muss, umso näher an seinem Arbeitsplatz sollten sie aufbewahrt werden. Die
Zugriffshäufigkeit entscheidet also über den Ablageort.

\textbf{(b)} Die drei Ablageorte:
\begin{itemize}[itemsep=0pt]
  \item \emph{Arbeitsplatzablage}: aktuell bearbeitete oder häufig benötigte Unterlagen direkt am Schreibtisch.
  \item \emph{Abteilungsablage}: von mehreren Mitarbeitern genutzte Unterlagen, gut erreichbar in der Abteilung (Aktenplan hilfreich).
  \item \emph{Zentralregistratur/Archiv}: Unterlagen aus verschiedenen Bereichen, selten benötigt oder mit Dauerwert.
\end{itemize}

\textbf{(c)} Die Zentralregistratur entlastet die Abteilungen von Schriftstücken, die
nicht mehr benötigt werden, aber noch aufbewahrt werden müssen (z.\,B.\ wegen
gesetzlicher Aufbewahrungspflicht). Außerdem werden hier Unterlagen mit Dauerwert
archiviert; ein Aktenplan sorgt für übersichtliche Ordnung.
\end{loesung}

\subsection*{Lösung Aufgabe~2}
\begin{loesung}
\textbf{(a)} Loseblattablage: $+$ Zeitgewinn, da nicht gelocht werden muss; $-$ Gefahr,
dass Unterlagen verloren gehen / längere Suchzeiten bei umfangreichen Akten.

\textbf{(b)} Geheftete Ablage: $+$ sichere Aufbewahrung und schnelles Wiederfinden;
$-$ höherer Zeitaufwand, da vor dem Abheften gelocht werden muss.

\textbf{(c)} \emph{Einzelakte}: Schriftstücke zu einem Vorgang/Fall, z.\,B.\ Personalakte
eines Mitarbeiters. \emph{Sammelakte}: Schriftstücke zu verschiedenen, aber
gleichartigen Vorgängen, z.\,B.\ Lieferscheine oder Rechnungen.

\textbf{(d)} Die Loseblattablage eignet sich für Schriftstücke, die im Arbeitsprozess
weitergegeben werden, nicht zu umfangreich sind und schnell zur Hand sein müssen --- hier
wäre das Lochen und Abheften unnötiger Aufwand.
\end{loesung}

\subsection*{Lösung Aufgabe~3}
\begin{loesung}
\textbf{(a)} \emph{Liegend}: Flachablage. \emph{Stehend}: Ordnerregistratur oder
Stehsammlerregistratur. \emph{Hängend}: Hängeregistratur oder Pendelregistratur.

\textbf{(b)} Beispiele für Schriftgutbehälter:
\begin{itemize}[itemsep=0pt]
  \item Flachablage: Aktendeckel, Schnellhefter, Jurismappe.
  \item Ordnerregistratur: Ordner. \quad Stehsammlerregistratur: Einstellmappen, Stehsammler, Kassetten.
  \item Hängeregistratur: Hängemappen, Hängetaschen, Hängeordner.
  \item Pendelregistratur: Pendelmappen, Pendeltaschen, Pendelhefter.
\end{itemize}

\textbf{(c)} Ordnerregistratur: $+$ sehr übersichtlich durch beschriftbare Ordnerrücken,
schneller Zugriff, geringe Kosten, sicheres Abheften; $-$ zeitaufwendig (Lochen und
Abheften) und bei unregelmäßigem Aktenanfall unflexibel (Platz muss frei bleiben).

\textbf{(d)} \emph{Hängeregistratur}: sehr übersichtlich, schneller Zugriff, aber
größerer Raumbedarf (nur bis Sichthöhe). \emph{Pendelregistratur}: bessere Ausnutzung
der Stellfläche (bis Griffhöhe), aber weniger übersichtlich (nur Schmalseite sichtbar)
und längere Zugriffszeit, da der ganze Behälter entnommen werden muss.
\end{loesung}

\subsection*{Lösung Aufgabe~4}
\begin{loesung}
\textbf{(a)} \textbf{Kaiser, Krähmer, Kuhn, Kurz, Larsson, Ludwig.}

\textbf{(b)} \textbf{Schade, Schäfer, Schmid, Schmidt, Schmitt, Schulze.} (Begründung:
Bei \emph{Schmid -- Schmidt} ist „Schmid” kürzer und steht voran; danach „Schmitt”, da
„dt” vor „tt” steht. Bei der Wörterbuch-Variante gilt ä = a, daher „Schade” vor
„Schäfer”.)

\textbf{(c)} Maßgeblich ist die Reihenfolge der Ordnungswörter; nach DIN 5007 entscheidet
die Buchstabenfolge des Alphabets (bei gleichem Anfang der nächste Buchstabe). Umlaute
werden in der Wörterbuch-Variante (Variante 1) wie der Grundbuchstabe behandelt (ä = a,
ö = o, ü = u), „ß” wie „ss”. (In der Namensverzeichnis-Variante 2 gilt ä = ae usw.)
\end{loesung}

\subsection*{Lösung Aufgabe~5}
\begin{loesung}
\begin{enumerate}[label=\arabic*., itemsep=0pt]
  \item \textbf{Alphabetisch} --- Kundennamen werden nach DIN 5007 alphabetisch geordnet (bei großen Beständen oft mit Kunden-Nr.\ kombiniert, also alphanumerisch).
  \item \textbf{Mnemotechnisch / alphanumerisch} --- die Kürzel ER/AR sind Gedächtnishilfen, die Rechnungsnummer dient der genauen Einordnung.
  \item \textbf{Chronologisch} --- Bilanzen und Steuerbescheide werden nach Jahren (zeitlicher Reihenfolge) abgelegt.
  \item \textbf{Geografisch} --- Ablage nach räumlichem Kriterium (Vertriebsgebiet), sinnvoll im Außendienst.
\end{enumerate}
\end{loesung}

\subsection*{Lösung Aufgabe~6}
\begin{loesung}
\textbf{(a)} In der Wiedervorlagemappe werden Aufgaben und Vorgänge abgelegt, die nicht
sofort bearbeitet werden (noch zu bearbeitende Schriftstücke, schwebende Vorgänge). Sie
vermeidet Stapel auf dem Schreibtisch und dient der Terminplanung. Aufbau: 12 Fächer für
die Monate oder 31 Fächer für die Monatstage; Trennblätter mit Sichtlöchern zeigen, ob
Schriftstücke enthalten sind.

\textbf{(b)} Drei weitere Ordnungshilfsmittel (Beispiele): \emph{Register-/Trennblätter}
(Übersicht im Ordner, alphabetisch/numerisch), \emph{Reiter} (Inhaltsschilder aus
Metall/Kunststoff auf dem Schriftgutbehälter), \emph{Leitkarten} (Unterteilung in Gruppen
bei Hängemappen/Karteikästen).

\textbf{(c)} Der Aktenplan (Informationsstrukturplan) ist der einheitliche Rahmen mit der
systematischen, logischen Ordnung für die Ablage. Bei der Dezimalklassifikation werden die
Hauptgruppen mit 0--9 nummeriert, Untergruppen und Stichworte durch weitere Ziffern
verfeinert (z.\,B.\ 5 Personal $\rightarrow$ 51 Personalverwaltung $\rightarrow$
514 Urlaubsanträge). Vorteil: einheitliche Regeln, leichte und sichere Einordnung.

\textbf{(d)} Die günstigste Registraturform spart zwar bei Material und Anschaffung, kann
aber höhere Arbeitszeit- und Raumkosten verursachen (z.\,B.\ langes Suchen, umständliche
Bearbeitung). Wirtschaftlich ist die Form, bei der Material-, Raum- und Zeitkosten
\emph{insgesamt} am niedrigsten sind --- nicht die mit dem billigsten Behälter.
\end{loesung}

\end{document}
